\section{动态抽屉原理}\label{sec:B}

在本附录中, 我们在分母为正的条件下, 
为在 \(\mathbf{Q}(x)\) 上线性无关的形式函数集合 \(f_1, \ldots, f_m \in \mathbf{Q}[\![x]\!]\) 
（具有形式 \eqref{eq:2.2.1} 并且满足 \(\varphi, \varphi^* f_1, \ldots, \varphi^* f_m \in \mathcal{M}(\overline{\mathbf{D}})\) 
在闭单位圆盘 \(\overline{\mathbf{D}}\) 的一个邻域内是联立亚纯的） 
给出了基本 holonomy 边界的简短新证明:
\begin{equation}\label{eq:B.0.1}
m \le \frac{2T(\varphi)}{\log|\varphi'(0)| - b_1 - \dots - b_r}.
\end{equation}
在此处,
\begin{equation}\label{eq:B.0.2}
T(\varphi) := \int_{\mathbf{T}} \log^{+} |\varphi| \, \mu_{\text{Haar}} + \sum_{\substack{\rho \in \mathbf{D} \\ \text{poles of } \varphi}} \log \frac{1}{\rho}
\end{equation}
是亚纯映射 \(\varphi\) 的 Nevanlinna 特征函数（亚纯极点按其重度计算）.

这基于 Perelli 以及 Zannier \cite{PZ84} 的一个动态抽屉原理想法, 
正如他们在该文献的引理 1 中所陈述的那样. 
它可以被视为第 7 节中 Bost 技术的更初等形式, 
它既是后者的引入也是一种替代方案, 
并且也指向我们的伴随论文 \cite{CDT24}, 
在其中这些想法得到了进一步的研究. 
我们根据第 2.12 节中的解剖将证明分为三个步骤.

\subsection{评估模}\label{sec:B.1}

假设我们拥有一组在 \(\mathbf{Q}(x)\) 上线性无关的、具有 \eqref{eq:2.2.1} 形式的函数 \(f_1(x), \dots, f_m(x) \in \mathbf{Q}[\![x]\!]\), 
使得 \(\varphi(z) \in \mathbf{C}[\![z]\!]\) 以及每个形式幂级数 \(f_i(\varphi(z)) \in \mathbf{C}[\![z]\!]\) 
都是闭复单位圆盘 \(|z| \le 1\) 邻域内的亚纯函数芽. 
（在用 \(\varphi(\rho z)\) 对于某些仍然满足 \(\rho|\varphi'(0)| > e^{b_1 + \dots + b_r}\) 的 \(\rho < 1\) 
代替 \(\varphi(z)\) 的前提下, 
拥有这个稍微大一些的圆盘并不损失一般性.） 
我们引入两个正整数参数 D 以及 T, 
并考虑集合:
\begin{equation}\label{eq:B.1.1}
\mathcal{I}_D(T) := \left\{ (Q_1, \dots, Q_m) \in \mathbf{Z}[x] : Q_i(x) = \sum_{j=0}^{D-1} c_{i,j} x^j, c_{i,j} \in [0, T) \cap \mathbf{Z} \right\},
\end{equation}
其基数为:
\begin{equation}\label{eq:B.1.2}
\#\mathcal{I}_D = T^{mD}.
\end{equation}

根据对 m 个形式幂级数 \(f_i(x) \in \mathbf{Q}[\![x]\!]\) 
所假设的 \(\mathbf{Q}(x)\)-线性无关性, 
\(\mathbf{Z}\)-模评估映射
\begin{equation}\label{eq:B.1.3}
\psi_D: \mathbf{Z}[x]_{\deg < D}^{\oplus m} \hookrightarrow \mathbf{Q}[\![x]\!], \qquad (Q_1, \dots, Q_m) \mapsto \sum_{i=1}^m Q_i(x) f_i(x) \in \mathbf{Q}[\![x]\!]
\end{equation}
是单射. 
因此, 该映射下的像
\[ \mathcal{O}_D := \psi_D(\mathcal{I}_D) \subset \mathbf{Q}[\![x]\!] \]
也具有基数:
\begin{equation}\label{eq:B.1.4}
\#\mathcal{O}_D = \#\mathcal{I}_D = T^{mD}.
\end{equation}

在秩为 mD 的自由 \(\mathbf{Z}\)-模 \eqref{eq:B.1.3} 中, 我们定义消逝过滤跳跃:
\begin{equation}\label{eq:B.1.5}
r_D^{(n)} := \dim_{\mathbf{Q}} \left\{ \left( \psi_D^{-1} \left( x^n \mathbf{Q} [\![x]\!] \right) \otimes \mathbf{Q} \right) / \left( \psi_D^{-1} \left( x^{n+1} \mathbf{Q} [\![x]\!] \right) \otimes \mathbf{Q} \right) \right\}.
\end{equation}
它们位于 \(\{0,1\}\) 中, 因为单射线性映射 \(\psi_D\) 诱导了一个线性单射:
\[ \left(\psi_D^{-1}\left(x^n\mathbf{Q}[\![x]\!]\right)\otimes\mathbf{Q}\right)/\left(\psi_D^{-1}\left(x^{n+1}\mathbf{Q}[\![x]\!]\right)\otimes\mathbf{Q}\right)\hookrightarrow x^n\mathbf{Q}[\![x]\!]/x^{n+1}\mathbf{Q}[\![x]\!]\cong\mathbf{Q}\cdot x^n \]
到一维 \(\mathbf{Q}\)-向量空间中. 
在另一方面, 我们有:
\begin{equation}\label{eq:B.1.6}
\sum_{n=0}^{\infty} r_D^{(n)} = \dim_{\mathbf{Q}} \psi_D^{-1} \mathbf{Q} [\![x]\!] = mD.
\end{equation}
因此, 存在一个大小为 mD 的、我们的辅助函数在 \(x = 0\) 处的可能消逝阶数集合:
\begin{equation}\label{eq:B.1.7}
\left\{ n \in \mathbf{N} : \exists (Q_1, \dots, Q_m) \in \mathbf{Z}[x]_{\deg < D}^{\oplus m}, \operatorname{ord}_{x=0} \left( \sum_{i=1}^m Q_i(x) f_i(x) \right) = n \right\}
\end{equation}
\[ = \left\{ 0 \le u(1) < u(2) < \dots < u(mD) \right\}. \]
它们的设定仅依赖于模 \eqref{eq:B.1.3} 
——换句话说, 仅依赖于 \(f_1, \ldots, f_m\) 以及参数 D, 
——而不依赖于参数 T, 该参数在后文中仍保持自由选择. 
（参数 T 将被选为依赖于过滤跳跃 \eqref{eq:B.1.7} 的任意足够大的整数.）

这完成了第 2.12 节的步骤 (i).

\subsection{抽屉原理}\label{sec:B.2}

对于步骤 (ii)，在临界条件 \(|\varphi'(0)| > e^{b_1 + \dots + b_r}\) 下，我们来测量辅助函数 \(F(x)\) 的 Taylor 级数递归地依赖于其初始系数串的趋势。

我们可以通过乘积 \(\prod_{p=1}^{mD} \gamma_p\) 来上估计输出基数 \(\#\mathcal{O}_D\)，其中 \(\gamma_p\) 是在共享一个公共先前系数串 \((\beta_0, \beta_1, \ldots, \beta_{u(p)-1})\) 的任意输出函数 \(F(x) = \sum_{k=0}^{\infty} \beta_k x^k\) 集合中、不同 \(x^{u(p)}\) 系数 \(\beta_{u(p)} \in \mathbf{Q}\) 的最大可能数量的上估计：
\[ \forall (\beta_0, \dots, \beta_{u(p)-1}) \in \mathbf{Q}^{u(p)}, \]
\[ \# \left\{ \beta \in \mathbf{Q} : \exists (Q_1, \dots, Q_m) \in \mathcal{I}_D, \sum_{i=1}^m Q_i(x) f_i(x) - \sum_{k=0}^{u(p)-1} \beta_k x^k = \beta x^{u(p)} + O(x^{u(p)+1}) \right\} \leq \gamma_p. \]

此时，利用整性条件 \eqref{eq:2.6.2} 可知，所有这些有理数 \(\beta \in \mathbf{Q}\) 实际上都属于由所需分母类型所给出的 \(\mathbf{Z}\)-模中：
\[ \beta \in \frac{1}{[1,\ldots,b_1 \cdot u(p)]\cdots[1,\ldots,b_r \cdot u(p)]}\mathbf{Z}. \]
因此，如果 \(A_p \in \mathbf{R}^{>0}\) 使得任意两个此类系数 \(\beta\) 的差的绝对值都在区间 \([-A_p, A_p]\) 内，那么我们可以取：
\begin{equation}\label{eq:B.2.1}
\gamma_p := 1 + 2A_p \cdot [1, \dots, b_1 u(p)] \cdots [1, \dots, b_r u(p)] = 1 + A_p \cdot e^{(b_1 + \dots + b_r)u(p) + o(u(p)) + O(1)}
\end{equation}
作为在给定 \((\beta_0, \dots, \beta_{u(p)-1})\) 的情况下 \(\beta_{u(p)}\) 所有可能输出值的总边界。这一步模仿了 Perelli 和 Zannier 的动态抽屉原理 \cite[第 2 节引理 1]{PZ84}（或译为动态盒原理）。

\subsection{对最低阶系数的丢番图分析}\label{sec:B.3}

对于 \(A_p\) 的边界，在解析上来自于使用多变量亚纯一致化映射 \(\varphi\)。考虑将 \(\varphi\) 表示为两个在闭单位圆盘 \(\overline{\mathbf{D}}\) 上收敛且满足 \(u(0) = 1\) 的幂级数 \(u\) 与 \(v\) 的商的任意表示：\(\varphi = v/u\)。令 \(h\) 为 \(\overline{\mathbf{D}}\) 上的一个满足 \(h(0) = 1\) 且使得对每个 \(i = 1, \ldots, m\)，\(hf_i\) 在 \(\overline{\mathbf{D}}\) 上全纯的收敛幂级数。为了书写简单，我们在下文中设 \(n := u(p)\)。任意两个上述输出函数 \(F_1(x)\) 和 \(F_2(x)\)，若其在 \(x = 0\) 处的 Taylor 级数一致直到 \(O(x^n)\)，且其相应的 \(x^n\) 系数分别为 \(\beta^{(1)}\) 和 \(\beta^{(2)}\)，那么它们将使得函数
\[ V(z) := h(z)u(z)^D \cdot (F_1(\varphi(z)) - F_2(\varphi(z))) \]
在闭单位圆盘 \(\overline{\mathbf{D}}\) 的一个邻域上是全纯的（收敛的），并且其领先阶项 \(\varphi'(0)^n(\beta^{(1)}-\beta^{(2)})z^n+O(z^{n+1})\) 可以根据 Cauchy 公式表示为在 \(|z|=1\) 上的围道积分：
\begin{equation}\label{eq:B.3.1}
\varphi'(0)^n(\beta^{(1)} - \beta^{(2)}) = \int_{\mathbf{T}} \frac{V(z)}{z^{n+1}} \,\mu_{\text{Haar}}.
\end{equation}

通过对被积函数的上确界进行估计，我们可以将 \(A_p\) 取为下式底部的上界：
\begin{equation}\label{eq:B.3.2}
\begin{aligned}
\left| \beta^{(1)} - \beta^{(2)} \right| &\le |\varphi'(0)|^{-n} \cdot \sup_{\mathbf{T}} |V| \\
&\le |\varphi'(0)|^{-u(p)} \cdot T \left( \sup_{\mathbf{T}} \max(|u|, |v|) \right)^{D} \cdot mD \cdot \sup_{\mathbf{T}} |h \cdot \varphi^* f_i| =: A_p,
\end{aligned}
\end{equation}
其中采用了 \(n := u(p)\)。

对于 \(\prod_{p=1}^{mD} \gamma_p\) 种可能输出，我们获得了总输出基数 \(T^{mD} = \#\mathcal{O}_D\) 的如下上估计：
\[
\begin{aligned}
\le \prod_{p=1}^{mD} \Big\{ &1 + T \exp\left(-\left(\log|\varphi'(0)| - b_1 - \dots - b_r + o(1)\right) u(p)\right) \\
&+ D \sup_{\mathbf{T}} \log \max(|u|, |v|) + \log D + O_{m,h}(1)\Big\}.
\end{aligned}
\]

此时，我们考虑上述不等式在 \(T \to \infty\) 下的渐近行为。更具体地，我们选择一个足够大的 \(T\)，使得该乘积的所有 \(mD\) 个因子都 \(\ge 2\)。利用当 \(x \ge 1\) 时的平凡不等式 \(1 + x \le 2x\)，并在等式两边消去共同出现的 \(T^{mD}\)，在取对数后，我们得到：
\[
\begin{aligned}
&-(\log |\varphi'(0)| - b_1 - \dots - b_r)(1 - o(1)) \sum_{p=1}^{mD} u(p) \\
&+ mD^2 \sup_{\mathbf{T}} \log \max(|u|, |v|) + O(D \log D) \ge -mD \log 2 - O_{m,h}(D).
\end{aligned}
\]

由于 \(0 \le u(1) < u(2) < \cdots < u(mD)\) 是非负整数的严格单增序列，我们有：
\begin{equation}\label{eq:B.3.3}
\sum_{p=1}^{mD} u(p) \ge \sum_{n=0}^{mD-1} n = \binom{mD}{2}.
\end{equation}

由此，我们导出：
\begin{equation}\label{eq:B.3.4}
(1-o(1))\binom{mD}{2}\left(\log|\varphi'(0)|-b_1-\dots-b_r\right) \le mD^2 \sup_{\mathbf{T}}\log\max(|u|,|v|)+O_{m,h}(D),
\end{equation}
在 \(D \to \infty\) 的渐近下，这最终过滤为如下算术 holonomy 边界：
\begin{equation}\label{eq:B.3.5}
m \le \frac{2\sup_{\mathbf{T}}\log\max(|u|,|v|)}{\log|\varphi'(0)| - b_1 - \dots - b_r}.
\end{equation}

该结论对于满足 \(u(0) = 1\) 的任何亚纯商表示 \(\varphi = v/u\) 均成立。Nevanlinna 的一个著名引理（参见 \cite[第 VII.1.4 节]{Nev70}），基于 Poisson--Jensen 公式中 canonical Blaschke 乘积以及对数分解 \(\log = \log^+ - \log^-\)，在开圆盘 \(\mathbf{D}\) 上构造了一个满足 \(u(0) = 1\)且满足 \(\sup_{\mathbf{D}} |u|\) 和 \(\sup_{\mathbf{D}} |v|\) 均由 \(\exp(T(\varphi))\) 所控制的商表示 \(\varphi = v/u\)。稍微放大（扩张）半径，对于任何 \(\varepsilon > 0\)，我们都可以获得一个商表示 \(\varphi = v/u\)，其定义在上述分析所需的闭圆盘 \(\mathbf{D}\) 的某个开邻域上，且满足 \(\sup_{\mathbf{T}} |u|\) and \(\sup_{\mathbf{T}} |v|\) 均由 \(\exp(T(\varphi) + \varepsilon)\) 所控制。这便完成了边界 \eqref{eq:B.0.1} 的证明。\(\square\)